\documentclass[dvipdfmx,11pt]{jsarticle}

\usepackage{amsmath, amssymb, amsthm}
\usepackage{ascmac}
\usepackage{color}
\usepackage{empheq}
\usepackage{mathrsfs}
\usepackage[margin=25truemm]{geometry}
\usepackage{setspace}

\usepackage{tikz}
\usetikzlibrary{calc, angles, quotes}




\newcommand{\Div}{\operatorname{div}}
\newcommand{\dsp}{\displaystyle}
\newcommand{\e}{\varepsilon}
\newcommand{\Ha}{\mathcal{H}}
\newcommand{\prob}{\vspace{2mm}\pro}
\newcommand{\redprob}{\vspace{2mm}\color{red}\pro\color{black}}
\newcommand{\blueprob}{\vspace{2mm}\color{blue}\pro\color{black}}
\renewcommand{\qedsymbol}{$\blacksquare$}
\newcommand{\R}{\mathbb{R}} 
\newcommand{\setmid}{\mathrel{}\middle|\mathrel{}}


\theoremstyle{definition}
\newtheorem{theorem}{Theorem}[section]
\newtheorem{lemma}[theorem]{Lemma}
\newtheorem{corollary}[theorem]{Corollary}
\newtheorem{Def}[theorem]{Definition}
\newtheorem{prop}[theorem]{Proposition}
\newtheorem{rem}[theorem]{Remark}
\newtheorem{ex}[theorem]{Example}
\newtheorem{assumption}{Assumption}
\newtheorem{pro}{問題}

\newtheorem*{theorem*}{Theorem}
\newtheorem*{lemma*}{Lemma}
\newtheorem*{corollary*}{Corollary}
\newtheorem*{prop*}{Proposition}
\newtheorem*{rem*}{Remark}
\newtheorem*{assumption*}{Assumption}
\newtheorem*{notation*}{Notation}

\DeclareMathOperator{\spt}{spt} 

\newcommand{\mres}{\mathbin{\vrule height 1.6ex depth 0pt width
0.13ex\vrule height 0.13ex depth 0pt width 1.3ex}}
\renewcommand{\labelenumi}{(\arabic{enumi})}

\makeatletter
\c@MaxMatrixCols=11
\makeatother

\numberwithin{equation}{section}

\setlength{\parindent}{0pt}

\title{第2回 関数の極限と連続性 演習問題}
\author{}
\date{}

\begin{document}

\maketitle

\vspace{-1.5cm}

\pro $f(x)=x^2$ が$\mathbb{R}$上で連続関数であることを定義から直接示せ.
	
\prob 以下の関数の連続性を調べよ.
	
(1) $f(x)=
\begin{cases}
		x & (x\in\mathbb{Q}) \\ 0 & (x\notin\mathbb{Q})
\end{cases}$

(2) $f(x)=
\begin{cases}
		\dsp \frac{1}{p} & (x=\dsp \frac{q}{p}\neq0 , p\in\mathbb{N} , q\in\mathbb{Z} , (p , q)=1) \\ 1 & (x=0) \\ 0 & (x\notin\mathbb{Q})
\end{cases}$

\prob 関数$f, g : \mathbb{R}\rightarrow\mathbb{R}$は$\dsp\lim_{x\rightarrow a}g(x)=b, \;\lim_{x\rightarrow b}f(x)=c$を満たす. このとき, $\dsp\lim_{x\rightarrow a}f(g(x))=c$が成立するか.

\prob 関数$f : (a, b)\rightarrow\mathbb{R}$は任意の$x_0\in[a,b]$に対して極限値$\dsp\lim_{x\rightarrow x_0}f(x)$が存在するとする. このとき, $f$が有界であることと, 任意の$x_0\in[a,b]$に対して極限値$\dsp\lim_{x\rightarrow x_0}f(x)$が有限であることは同値であることを示せ.

\prob $f(x) , g(x)$ を実数上で定義された連続関数とし, すべての有理数 $a$ について $f(a) = g(a)$ が成立するとする. このときすべての実数 $x$ について $f(x)=g(x)$ であることを示せ.
	
\prob 区間 $I$ で定義される関数 $f , g$ に対し, $f と g$ が $x_0\in I$ で連続かつ $f(x_0)>g(x_0)$ ならば, ある $\delta >0$ が存在して, 任意の $x\in (x_0-\delta , x_0+\delta)\cap I$ に対し $f(x)>g(x)$ が成立することを示せ.

\prob 関数$f,g : [a,b]\rightarrow(0, \infty)$は連続で, 任意の$x\in[a,b]$に対して$g(x)<f(x)$が成立するとする. このときある$\lambda>1$が存在して, $[a,b]$上で$f\geq\lambda g$が成立することを示せ.

\prob 関数 $f : (a,\infty)\rightarrow \mathbb{R}$ が $\dsp\lim_{x\rightarrow\infty}xf(x)=L\in\mathbb{R}$ を満たすならば, $\dsp\lim_{x\rightarrow\infty}f(x)=0$ 
であることを示せ.

\prob 区間上定義された単射連続関数は狭義単調増加であることを示せ.

\prob 連続関数 $f : \mathbb{R}\rightarrow \mathbb{R}$ に対し, ある実数 $\alpha_+ , \alpha_-$ が存在し, $\dsp \lim_{x\rightarrow\pm\infty}f(x)=\alpha_\pm$ となるならば $f$ は有界であることを示せ. また, $\alpha_+=\alpha_-$ のとき $f$ は最大値または最小値をもつことを示せ.

\prob 関数$f, g :\mathbb{R}\rightarrow\mathbb{R}$は周期関数で, $\dsp\lim_{x\rightarrow\infty}(f(x)-g(x))=0$を満たすとする.

(1) $f$と$g$は同じ周期を持つことを示せ.

(2) $f=g$を示せ.

\prob $\mathbb{R}$ 上定義された連続な周期関数は有界か.

\prob $f(x)$ は区間 $[a , b]$ で連続とするとき, $\dsp\int_{a}^{b}|f(x)|dx=0$ ならば, 任意の $x\in [a , b]$ に対して $f(x)=0$ であることを示せ. また, $f$ が区間 $[a , b]$ で積分可能であるが連続でないときは成立するか.
		
\prob 連続関数の区間の像は区間か1点集合であることを示せ.

\prob 連続関数 $f : \mathbb{R}\rightarrow \mathbb{R}$は開区間を開区間に写すとする. このとき$f$は単調であることを示せ.
		
\prob 任意の実係数奇数次多項式は少なくとも1つの実根をもつことを示せ.
		
\prob $f : [0 , 1]\rightarrow [0 , 1]$ を連続関数とする. このとき不動点, すなわち $f(x)=x$ となる $x\in [0 , 1]$ が存在することを示せ.

\prob 実数上の連続関数で, 任意の $x , y\in\mathbb{R}$ について以下の条件を満たすものをそれぞれ答えよ.
	
(1) $f(x+y)=f(x)+f(y)$
	
(2) $f(x)>0$ かつ $f(x+y)=f(x)f(y)$
	
(3) $f(x)>0$ かつ $f(xy)=f(x)+f(y)$
	
(4) $f(x)>0$ かつ $f(xy)=f(x)f(y)$

\prob 実数上の関数で, 以下のものを満たすものは定数関数であることを示せ.
	
(1) $f(2x)=f(x)$ かつ $x=0$ で連続.
	
(2) $f(x^2)=f(x)$ かつ $x=0$ および $x=1$ で連続.

\prob 
(1) $f(x)=\sin x$ は $(-\infty , \infty)$ で一様連続であることを示せ.

(2) $f(x)=x^2$ は $(0 , 1)$ で一様連続であるが, $[1 , \infty)$ で一様連続でないことを示せ.
	
(3) $\dsp f(x)=\frac{1}{x}$ は $(0 , 1)$ で一様連続でないが, $[1 , \infty)$ で一様連続であることを示せ.
	
(4) $f(x)=x^3-2$ は $(-\infty , \infty)$ で一様連続でないことを示せ.

\end{document}